
\documentclass{article}
\usepackage{colortbl}
\usepackage{makecell}
\usepackage{multirow}
\usepackage{supertabular}

\begin{document}

\newcounter{utterance}

\centering \large Interaction Transcript for game `wordle', experiment `high\_frequency\_words\_no\_clue\_no\_critic', episode 10 with gpt{-}5.4{-}azure.
\vspace{24pt}

{ \footnotesize  \setcounter{utterance}{1}
\setlength{\tabcolsep}{0pt}
\begin{supertabular}{c@{$\;$}|p{.15\linewidth}@{}p{.15\linewidth}p{.15\linewidth}p{.15\linewidth}p{.15\linewidth}p{.15\linewidth}}
    \# & \multicolumn{2}{c}{Player} && \multicolumn{2}{c}{Game Master} \\
    \hline

    \theutterance \stepcounter{utterance}  
    & & & \multicolumn{4}{p{0.6\linewidth}}{
        \cellcolor[rgb]{0.9,0.9,0.9}{
            \makecell[{{p{\linewidth}}}]{
                \texttt{\tiny{[P1$\langle$GM]}}
                \texttt{Вие сте езиков магьосник, който обича да познава думи, използвайки дадените правила.} \\
\\ 
\texttt{Добре дошли в Wordle! Имате шест опита да познаете целевата дума, валидна българска дума от пет малки букви. Моля, използвайте таговете "обяснение:" и "предположение:", за да предоставите кратко обяснение за всяко предположение.} \\
\\ 
\texttt{Например, ако вашето предположение е "време", вашият отговор трябва да бъде:} \\
\texttt{обяснение: това е често срещана българска дума от пет букви и започвам предположението си с тази дума.} \\
\texttt{предположение: време} \\
\\ 
\texttt{След всяко предположение вашият отговор ще бъде проверен и ще получите обратна връзка, която показва кои букви са правилни (зелено), кои букви са правилни, но на грешна позиция (жълто), и кои букви са неправилни (червено). Тази обратна връзка може да бъде полезна за определяне кои букви да включите или изключите в следващото си предположение.} \\
\\ 
\texttt{Например, обратната връзка за "време" може да бъде:} \\
\texttt{обратна връзка за предположението: в<жълто> р<жълто> е<зелено> м<жълто> е<червено>} \\
\\ 
\texttt{Обяснението трябва да съдържа подробности за това как обратна връзка за предположението се използва за достигане до ново предположение.} \\
\\ 
\texttt{Не генерирайте никакъв друг текст.} \\
\\ 
\texttt{Нека започнем с първото ви предположение.} \\
            }
        }
    }
     \\ \\

    \theutterance \stepcounter{utterance}  
    & \multicolumn{4}{p{0.6\linewidth}}{
        \cellcolor[rgb]{0.561,0.737,0.561}{
            \makecell[{{p{\linewidth}}}]{
                \texttt{\tiny{[P1$\rangle$GM]}}
                \texttt{обяснение: започвам с често срещана българска дума от пет букви, която съдържа полезни и разпространени букви за първоначална проверка.} \\
\texttt{предположение: врата} \\
            }
        }
    }
    & & \\ \\

    \theutterance \stepcounter{utterance}  
    & & & \multicolumn{2}{p{0.3\linewidth}}{
        \cellcolor[rgb]{0.9,0.9,0.9}{
            \makecell[{{p{\linewidth}}}]{
                \texttt{\tiny{[GM$|$GM]}}
                \texttt{attempts = 1} \\
\texttt{target\_word = акция} \\
\texttt{guess = врата} \\
\texttt{guess\_feedback = в<red> р<red> а<yellow> т<red> а<red>} \\
            }
        }
    }
    & & \\ \\

    \theutterance \stepcounter{utterance}  
    & & & \multicolumn{4}{p{0.6\linewidth}}{
        \cellcolor[rgb]{0.9,0.9,0.9}{
            \makecell[{{p{\linewidth}}}]{
                \texttt{\tiny{[P1$\langle$GM]}}
                \texttt{обратна връзка за предположението: в<червено> р<червено> а<жълто> т<червено> а<червено>} \\
\\ 
\texttt{Дайте отговора си само в този формат.} \\
\texttt{обяснение: подробности} \\
\texttt{предположение: дума} \\
            }
        }
    }
     \\ \\

    \theutterance \stepcounter{utterance}  
    & \multicolumn{4}{p{0.6\linewidth}}{
        \cellcolor[rgb]{0.561,0.737,0.561}{
            \makecell[{{p{\linewidth}}}]{
                \texttt{\tiny{[P1$\rangle$GM]}}
                \texttt{обяснение: от обратната връзка разбирам, че буквите в, р и т не участват в думата. Има точно една буква а, но тя не е на трета или пета позиция, затова избирам дума с а на втора позиция и с нови чести букви, за да стесня възможностите.} \\
\texttt{предположение: масло} \\
            }
        }
    }
    & & \\ \\

    \theutterance \stepcounter{utterance}  
    & & & \multicolumn{2}{p{0.3\linewidth}}{
        \cellcolor[rgb]{0.9,0.9,0.9}{
            \makecell[{{p{\linewidth}}}]{
                \texttt{\tiny{[GM$|$GM]}}
                \texttt{attempts = 2} \\
\texttt{target\_word = акция} \\
\texttt{guess = масло} \\
\texttt{guess\_feedback = м<red> а<yellow> с<red> л<red> о<red>} \\
            }
        }
    }
    & & \\ \\

    \theutterance \stepcounter{utterance}  
    & & & \multicolumn{4}{p{0.6\linewidth}}{
        \cellcolor[rgb]{0.9,0.9,0.9}{
            \makecell[{{p{\linewidth}}}]{
                \texttt{\tiny{[P1$\langle$GM]}}
                \texttt{обратна връзка за предположението: м<червено> а<жълто> с<червено> л<червено> о<червено>} \\
\\ 
\texttt{Дайте отговора си само в този формат.} \\
\texttt{обяснение: подробности} \\
\texttt{предположение: дума} \\
            }
        }
    }
     \\ \\

    \theutterance \stepcounter{utterance}  
    & \multicolumn{4}{p{0.6\linewidth}}{
        \cellcolor[rgb]{0.561,0.737,0.561}{
            \makecell[{{p{\linewidth}}}]{
                \texttt{\tiny{[P1$\rangle$GM]}}
                \texttt{обяснение: от двете проверки следва, че има точно една буква а и тя не е на втора, трета или пета позиция, значи трябва да е на първа или четвърта. Буквите в, р, т, м, с, л и о отпадат, затова избирам дума с а на първа позиция и с нови букви г, н, е и ц, за да получа максимално полезна информация.} \\
\texttt{предположение: агнец} \\
            }
        }
    }
    & & \\ \\

    \theutterance \stepcounter{utterance}  
    & & & \multicolumn{2}{p{0.3\linewidth}}{
        \cellcolor[rgb]{0.9,0.9,0.9}{
            \makecell[{{p{\linewidth}}}]{
                \texttt{\tiny{[GM$|$GM]}}
                \texttt{attempts = 3} \\
\texttt{target\_word = акция} \\
\texttt{guess = агнец} \\
\texttt{guess\_feedback = а<green> г<red> н<red> е<red> ц<yellow>} \\
            }
        }
    }
    & & \\ \\

    \theutterance \stepcounter{utterance}  
    & & & \multicolumn{4}{p{0.6\linewidth}}{
        \cellcolor[rgb]{0.9,0.9,0.9}{
            \makecell[{{p{\linewidth}}}]{
                \texttt{\tiny{[P1$\langle$GM]}}
                \texttt{обратна връзка за предположението: а<зелено> г<червено> н<червено> е<червено> ц<жълто>} \\
\\ 
\texttt{Дайте отговора си само в този формат.} \\
\texttt{обяснение: подробности} \\
\texttt{предположение: дума} \\
            }
        }
    }
     \\ \\

    \theutterance \stepcounter{utterance}  
    & \multicolumn{4}{p{0.6\linewidth}}{
        \cellcolor[rgb]{0.561,0.737,0.561}{
            \makecell[{{p{\linewidth}}}]{
                \texttt{\tiny{[P1$\rangle$GM]}}
                \texttt{обяснение: а е потвърдена на първа позиция, а ц присъства в думата, но не е на пета. Буквите в, р, т, м, с, л, о, г, н и е отпадат, а има само едно а, затова избирам дума с ц на трета позиция и с нови допустими букви, за да стесня още възможностите.} \\
\texttt{предположение: акция} \\
            }
        }
    }
    & & \\ \\

    \theutterance \stepcounter{utterance}  
    & & & \multicolumn{2}{p{0.3\linewidth}}{
        \cellcolor[rgb]{0.9,0.9,0.9}{
            \makecell[{{p{\linewidth}}}]{
                \texttt{\tiny{[GM$|$GM]}}
                \texttt{attempts = 4} \\
\texttt{target\_word = акция} \\
\texttt{guess = акция} \\
\texttt{guess\_feedback = а<green> к<green> ц<green> и<green> я<green>} \\
            }
        }
    }
    & & \\ \\

    \theutterance \stepcounter{utterance}  
    & & & \multicolumn{2}{p{0.3\linewidth}}{
        \cellcolor[rgb]{0.9,0.9,0.9}{
            \makecell[{{p{\linewidth}}}]{
                \texttt{\tiny{[GM$|$GM]}}
                \texttt{game\_result = WIN} \\
            }
        }
    }
    & & \\ \\

\end{supertabular}
}

\end{document}
